% GENERATED FILE. Edit canonical lessons, then run npm run generate:books.
% canonical-source-hash: fnv1a64:7ad5df78006f6cc2
% canonical-lessons: LA-C128-placenta, LA-C128-cervisia, LA-C128-sucus, LA-C128-rosa, LA-C128-lilium

\chapter{Flat cake, Beer, Juice, Rose, Lily}
\label{ch:la-flat-cake-beer-juice-rose-lily}
\hlchaptermodality{128}

\begin{chapteropening}
\textbf{By the end of this chapter:} \emph{I can say a flat cake, beer, juice, a rose and a lily.}

Complete the last lesson of chapter 128: I can say a flat cake, beer, juice, a rose and a lily.
\end{chapteropening}

\section[placenta, placentae]{placenta, placentae — a flat cake}

\label{lesson:LA-C128-placenta}



\begin{quote}
Before the new one: say the Latin for a chicken, then the Latin for pepper.
\end{quote}

\subsection*{You'll want to know: placenta, placentae}
\textbf{placenta, placentae} — \textquotedblleft{}a flat cake\textquotedblright{}.

\begin{cousinweb}
\textbf{placenta} was borrowed from Greek \emph{plakous}, \emph{plakountos}, \textquotedblleft{}a flat cake\textquotedblright{} (Lewis \& Short). Anatomists later gave the name to a flat organ of the womb.
\end{cousinweb}

\begin{grammarlens}[title={What you've built}]
One more word for this chapter.
\end{grammarlens}

\subsection*{Your turn}
\begin{itemize}
\raggedright
  \item \emph{Say it:} \emph{placenta}
  \item \emph{Say it:} \emph{placenta}, once more
  \item \emph{Say it:} say \emph{placenta}
  \item \emph{Recall:} say the Latin for a chicken, then the Latin for pepper, then say \emph{placenta} again
\end{itemize}

\begin{tcolorbox}[breakable,colback=teal!4,colframe=teal!35!black,title={Before you move on}]
What does placenta, placentae mean? (\textquotedblleft{}A flat cake\textquotedblright{}.) Say it once more.
\end{tcolorbox}
\section[cervīsia, cervīsiae]{cervīsia, cervīsiae — beer}

\label{lesson:LA-C128-cervisia}



\begin{quote}
Before the new one: say the Latin for pepper, then the Latin for a flat cake.
\end{quote}

\subsection*{You'll want to know: cervīsia, cervīsiae}
\textbf{cervīsia, cervīsiae} — \textquotedblleft{}beer\textquotedblright{}.

\begin{cousinweb}
\textbf{cervīsia} came from Gaulish, the Celtic language of Gaul: Lewis \& Short mark \emph{cervīsia} as Gallic, a Celtic word for beer.
\end{cousinweb}

\begin{grammarlens}[title={What you've built}]
One more word for this chapter.
\end{grammarlens}

\subsection*{Your turn}
\begin{itemize}
\raggedright
  \item \emph{Say it:} \emph{cervīsia}
  \item \emph{Say it:} \emph{cervīsia}, once more
  \item \emph{Say it:} say \emph{cervīsia}
  \item \emph{Recall:} say the Latin for pepper, then the Latin for a flat cake, then say \emph{cervīsia} again
\end{itemize}

\begin{tcolorbox}[breakable,colback=teal!4,colframe=teal!35!black,title={Before you move on}]
What does cervīsia, cervīsiae mean? (\textquotedblleft{}Beer\textquotedblright{}.) Say it once more.
\end{tcolorbox}
\section[sūcus, sūcī]{sūcus, sūcī — juice}

\label{lesson:LA-C128-sucus}



\begin{quote}
Before the new one: say the Latin for a flat cake, then the Latin for beer.
\end{quote}

\subsection*{You'll want to know: sūcus, sūcī}
\textbf{sūcus, sūcī} — \textquotedblleft{}juice\textquotedblright{}.

\begin{cousinweb}
\textbf{sūcus} belongs with \emph{sūgō}, \textquotedblleft{}to suck\textquotedblright{} (Lewis \& Short), and both are usually traced to Proto-Indo-European \emph{*sug-}, \textquotedblleft{}to suck\textquotedblright{} (EtymDB 2.1). English \emph{suck} is a cousin.
\end{cousinweb}

\begin{grammarlens}[title={What you've built}]
One more word for this chapter.
\end{grammarlens}

\subsection*{Your turn}
\begin{itemize}
\raggedright
  \item \emph{Say it:} \emph{sūcus}
  \item \emph{Say it:} \emph{sūcus}, once more
  \item \emph{Say it:} say \emph{sūcus}
  \item \emph{Recall:} say the Latin for a flat cake, then the Latin for beer, then say \emph{sūcus} again
\end{itemize}

\begin{tcolorbox}[breakable,colback=teal!4,colframe=teal!35!black,title={Before you move on}]
What does sūcus, sūcī mean? (\textquotedblleft{}Juice\textquotedblright{}.) Say it once more.
\end{tcolorbox}
\section[rosa, rosae]{rosa, rosae — a rose}

\label{lesson:LA-C128-rosa}



\begin{quote}
Before the new one: say the Latin for beer, then the Latin for juice.

Say \emph{amō}, I love, and then \emph{amīcum amō}.
\end{quote}

\subsection*{You'll want to know: rosa, rosae}
\textbf{rosa, rosae} — \textquotedblleft{}a rose\textquotedblright{}.

\begin{cousinweb}
\textbf{rosa} is related to Greek \emph{rhodon}, \textquotedblleft{}rose\textquotedblright{} (Lewis \& Short), but it is not a direct borrowing: it probably reached Latin through another language of Italy. The word is often traced further back to Iranian.
\end{cousinweb}

\begin{grammarlens}[title={What you've built}]
One more word for this chapter.
\end{grammarlens}

\subsection*{Your turn}
\begin{itemize}
\raggedright
  \item \emph{Say it:} \emph{rosa}
  \item \emph{Say it:} \emph{rosa}, once more
  \item \emph{Say it:} say \emph{rosa}
  \item \emph{Recall:} say the Latin for beer, then the Latin for juice, then say \emph{rosa} again
\end{itemize}

\begin{tcolorbox}[breakable,colback=teal!4,colframe=teal!35!black,title={Before you move on}]
What does rosa, rosae mean? (\textquotedblleft{}A rose\textquotedblright{}.) Say it once more.
\end{tcolorbox}
\section[līlium, līliī]{līlium, līliī — a lily}

\label{lesson:LA-C128-lilium}



\begin{quote}
Before the new one: say the Latin for juice, then the Latin for a rose.
\end{quote}

\subsection*{You'll want to know: līlium, līliī}
\textbf{līlium, līliī} — \textquotedblleft{}a lily\textquotedblright{}.

\begin{cousinweb}
\textbf{līlium} was borrowed from Greek \emph{leirion}, \textquotedblleft{}lily\textquotedblright{} (Lewis \& Short), probably through another language; the word may go back to Egyptian.
\end{cousinweb}

\begin{grammarlens}[title={What you've built}]
One more word for this chapter.
\end{grammarlens}

\subsection*{Your turn}
\begin{itemize}
\raggedright
  \item \emph{Say it:} \emph{līlium}
  \item \emph{Say it:} \emph{līlium}, once more
  \item \emph{Say it:} say \emph{līlium}
  \item \emph{Recall:} say the Latin for juice, then the Latin for a rose, then say \emph{līlium} again
\end{itemize}

\begin{tcolorbox}[breakable,colback=teal!4,colframe=teal!35!black,title={Before you move on}]
What does līlium, līliī mean? (\textquotedblleft{}A lily\textquotedblright{}.) Say it once more.
\end{tcolorbox}
